% ============================================================
% The Spectral Repulsion
% ============================================================

\documentclass[12pt]{amsart}

\usepackage{amsmath,amssymb,amsthm}
\usepackage{booktabs}
\usepackage{microtype}
\usepackage{url}

\newtheorem{theorem}{Theorem}
\newtheorem{lemma}[theorem]{Lemma}
\newtheorem{proposition}[theorem]{Proposition}
\newtheorem{corollary}[theorem]{Corollary}
\theoremstyle{definition}
\newtheorem{definition}[theorem]{Definition}
\newtheorem{observation}[theorem]{Observation}
\theoremstyle{remark}
\newtheorem{remark}[theorem]{Remark}

\title{The Spectral Repulsion}

\author{Alexander S.\ Petty}
\email{alexander.petty@gmail.com}
\date{October 2025 (revised April 2026)}
\subjclass[2020]{11A63, 11N05, 11M06}

\begin{document}
\begin{abstract}
The collision transform decomposes the centered
collision deviation $S^{\circ}$ over odd Dirichlet
characters modulo $b^{\ell+1}$. The squared
coefficients $|\hat{S}^{\circ}(\chi)|^2$ and the
squared prime character sums $|P(s, \chi)|^2$ define
two probability distributions on the same character
space. This paper measures their overlap.

The overlap decays like $\rho / n$, where $n$ is the
number of odd characters and $\rho \approx 0.6$ is
stable across all prime bases tested ($b = 3$ to $37$,
$n = 3$ to $666$). Equivalently, the scaled overlap
$n\Omega$ stays near $0.6$ across the tested bases.
Under the random-relabeling null model, where one
spectrum is uniformly permuted relative to the other,
the expected overlap is exactly $1/n$. Thus the
measured overlap is about $60\%$ of that null-model
baseline, and the missing $40\%$ is the spectral
repulsion.

The repulsion has consequences. The collision-weighted
Goldbach correlation $\sum_{p+q=N} S^{\circ}(p)\,
S^{\circ}(q)$ has near-zero mean in the tested range.
The reflection identity creates a further
obstruction: if $S(a) \ge 0$ then $S(m{-}a) \le -1$,
so both primes in a Goldbach pair cannot lie in the
same collision half-group.
\end{abstract}
\maketitle

% ============================================================
\section{Introduction}

The preceding papers recorded two computational
patterns: the anti-correlation between
$|\hat{S}^{\circ}(\chi)|$ and $|P(s, \chi)|$~\cite{paper19},
and the double transversality at the level of finite
base experiments~\cite{paper20}. Those notes do not
prove a general mechanism; they isolate a repeatable
numerical pattern.

This paper quantifies the avoidance. The two squared
distributions $|\hat{S}^{\circ}|^2$ and $|P|^2$ on
the odd characters have a measurable overlap that
decays with the group size. The scaled overlap
$n\Omega$ is the statistic studied here.

% ============================================================
\section{The Overlap}

\begin{definition}
For a fixed base $b$ at lag $\ell = 1$, let
$n = \phi(b^2)/2$ be the number of odd characters
modulo $b^2$. Define the normalized distributions
\[
p_c = \frac{|\hat{S}^{\circ}(\chi_c)|^2}
{\sum_j |\hat{S}^{\circ}(\chi_j)|^2},
\qquad
q_c = \frac{|P(s, \chi_c)|^2}
{\sum_j |P(s, \chi_j)|^2},
\]
where the sums run over odd characters. The
\emph{spectral overlap} is
\[
\Omega(b, s) = \sum_{c=1}^{n} p_c\, q_c.
\]
\end{definition}

The quantity $\Omega$ is the raw overlap of the two
normalized spectra. The comparison benchmark used in
this paper is the \emph{random-relabeling null model}:
fix the multisets of weights $\{p_c\}$ and $\{q_c\}$,
and choose a uniform random permutation $\pi$ of the
channels. Then the relabeled overlap is
\[
\Omega_\pi=\sum_{c=1}^{n} p_c\, q_{\pi(c)}.
\]
Since $\mathbb{E}[q_{\pi(c)}]=\frac1n\sum_j q_j=1/n$,
one has
\[
\mathbb{E}[\Omega_\pi]=\sum_{c=1}^{n} p_c\cdot\frac1n
=\frac1n.
\]
Thus $1/n$ is the exact expected overlap under random
relabeling, regardless of how nonuniform $p$ and $q$
are.

% ============================================================
\section{The Repulsion Ratio}

\begin{table}[h]
\centering
\begin{tabular}{rrrrrr}
\toprule
$b$ & $n$ & $\Omega$ & $1/n$ & $n\Omega$ &
$\Omega/(1/n)$ \\
\midrule
$3$ & $3$ & $0.1931$ & $0.3333$ & $0.58$ & $0.58$ \\
$5$ & $10$ & $0.0766$ & $0.1000$ & $0.77$ & $0.77$ \\
$7$ & $21$ & $0.0257$ & $0.0476$ & $0.54$ & $0.54$ \\
$11$ & $55$ & $0.0108$ & $0.0182$ & $0.60$ & $0.60$ \\
$13$ & $78$ & $0.0078$ & $0.0128$ & $0.61$ & $0.61$ \\
$17$ & $136$ & $0.0041$ & $0.0074$ & $0.56$ & $0.56$ \\
$19$ & $171$ & $0.0046$ & $0.0058$ & $0.79$ & $0.79$ \\
$23$ & $253$ & $0.0024$ & $0.0040$ & $0.60$ & $0.60$ \\
$29$ & $406$ & $0.0015$ & $0.0025$ & $0.63$ & $0.63$ \\
$31$ & $465$ & $0.0013$ & $0.0022$ & $0.63$ & $0.63$ \\
$37$ & $666$ & $0.0009$ & $0.0015$ & $0.58$ & $0.58$ \\
\bottomrule
\end{tabular}
\caption{Spectral overlap $\Omega$ at $s = 0.5$
across prime bases. The scaled overlap $n\Omega$ is
stable near $0.6$. Thus the measured overlap is about
$60\%$ of the random-relabeling benchmark $1/n$.}
\label{tab:overlap}
\end{table}

\begin{definition}
The \emph{repulsion ratio} is
$\rho(b,s) = n\,\Omega(b,s)$, the scaled overlap
at base~$b$ and point~$s$. If the limit
$\rho(s) = \lim_{n \to \infty} \rho(b,s)$ exists,
it is the asymptotic repulsion ratio.
\end{definition}

In the tested bases, $\rho \approx 0.6$ at $s = 0.5$,
with no visible dependence on~$b$. If $\rho = 1$, the distributions would be
as overlapped as the random-relabeling benchmark. If $\rho = 0$,
they would be perfectly orthogonal. The observed value
$\rho \approx 0.6$ indicates a $40\%$ deficit relative
to the random-relabeling benchmark: in the tested
bases, the collision spectrum under-overlaps the prime
spectrum.

\begin{remark}
The value $\rho \approx 0.6$ is close to
$1/\varphi = (\sqrt{5}-1)/2 \approx 0.618$. Whether
this is a coincidence or a structural connection to
the golden ratio (which selects the prime $3$ in the
alignment theory~\cite{paperA}) is an open question.
\end{remark}

% ============================================================
\section{Effective Support}

The Shannon entropy $H = -\sum p_c \log p_c$ measures
the spread of a distribution. The effective support is
$e^H$: the number of characters that carry most of the
distribution's weight.

\begin{table}[h]
\centering
\begin{tabular}{rrrrrr}
\toprule
$b$ & $n$ & eff($\hat{S}$) & eff($P$) &
product & product/$n$ \\
\midrule
$3$ & $3$ & $2.0$ & $2.5$ & $5.1$ & $1.7$ \\
$7$ & $21$ & $7.0$ & $10.7$ & $74.6$ & $3.6$ \\
$13$ & $78$ & $18.0$ & $52.0$ & $933$ & $12.0$ \\
$19$ & $171$ & $33.2$ & $112$ & $3{,}716$ & $21.7$ \\
$29$ & $406$ & $64.0$ & $295$ & $18{,}877$ & $46.5$ \\
$37$ & $666$ & --- & --- & --- & --- \\
\bottomrule
\end{tabular}
\caption{Effective support of the collision spectrum
and the prime spectrum. The collision spectrum
concentrates; the prime spectrum spreads. Their
product grows faster than $n$.}
\label{tab:effective}
\end{table}

The collision invariant concentrates its Fourier energy
on a small number of characters, with effective support
growing much more slowly than the ambient dimension in
the displayed data. The prime distribution spreads its
energy across most characters in the same data. The two
spectra therefore occupy visibly different regions of
character space in the tested bases.

% ============================================================
\section{The Goldbach Null}

The spectral repulsion has implications for the
collision invariant's interaction with additive
prime problems.

\begin{definition}
The \emph{collision-weighted Goldbach correlation} is
\[
H(N) = \sum_{\substack{p + q = N \\ p, q \text{ prime}}}
S^{\circ}(p)\, S^{\circ}(q).
\]
\end{definition}

\begin{observation}[Goldbach null]\leavevmode\newline
The normalized correlation
$H(N)/R(N)$ (where $R(N)$ is the Goldbach count)
averages to approximately zero across all tested
bases.
\end{observation}

\begin{table}[h]
\centering
\begin{tabular}{rrr}
\toprule
$b$ & mean $H/R$ & std $H/R$ \\
\midrule
$3$ & $+0.000$ & $0.63$ \\
$5$ & $+0.001$ & $0.89$ \\
$7$ & $+0.002$ & $1.20$ \\
$10$ & $+0.002$ & $1.55$ \\
$13$ & $+0.004$ & $1.75$ \\
\bottomrule
\end{tabular}
\caption{The collision-weighted Goldbach correlation
averages to zero. Even $N$ from $1{,}000$ to
$100{,}000$.}
\label{tab:goldbach}
\end{table}

The null result is consistent with the same separation
seen in the overlap statistic: the collision-weighted
Goldbach correlation does not show a sustained mean in
the tested range.

% ============================================================
\section{The Reflection Obstruction}

The reflection identity creates a further constraint
on collision-class Goldbach problems.

\begin{proposition}\label{prop:obstruction}
If $S(a) \ge 0$, then $S(m{-}a) \le -1$. Therefore,
for any Goldbach pair $p + q = N$ with
$p \equiv a \pmod{m}$ and $q \equiv N{-}a \pmod{m}$,
the two primes cannot both have non-negative collision
deviation when $N \equiv 0 \pmod{m}$.
\end{proposition}

\begin{proof}
By the reflection identity~\cite{paperA},
$S(a) + S(m{-}a) = -1$. If $S(a) \ge 0$, then
$S(m{-}a) = -1 - S(a) \le -1$.
When $N \equiv 0 \pmod{m}$, the residue of $q$ is
$m - a \pmod{m}$, so the two deviations have
opposite signs.
\end{proof}

\begin{remark}
For general $N$, the residue of $q$ is
$(N - a) \bmod m \ne m - a$ unless $N \equiv 0$.
The obstruction is partial, not total: it affects
$1/m$ of all even $N$. But it demonstrates that the
bilateral symmetry of the collision invariant is
structurally incompatible with same-class Goldbach
restrictions.
\end{remark}

% ============================================================
\section{Remarks}

The spectral repulsion quantifies the central
observation of this program: the collision invariant
and the prime distribution occupy different parts of
character space. The repulsion ratio $\rho \approx 0.6$
is stable across all bases tested and is the main
numerical invariant recorded by this note.

The collision invariant does not see the $L$-function
zeros in the same way the prime spectrum does
(anti-correlation, observed in the earlier note). Its
collision-weighted Goldbach correlation has near-zero
mean in the tested range. And the double-transversality
note suggests the same kind of misalignment across
bases. Together these observations point toward a
base-specific geometry that is not well captured by
standard additive or multiplicative averages alone.

The repulsion ratio $\rho \approx 0.6$ is the
quantitative signature of this separation. It
measures the degree to which the digit function's
geometry is transverse to the prime distribution's
analytic structure in the tested family. That this
ratio is stable across bases and close to a familiar
constant is suggestive, but no theorem identifying its
source is proved here.

% ============================================================
\begin{thebibliography}{99}

\bibitem{paperA}
A.~S.~Petty,
The collision invariant,
preprint, 2026.

\bibitem{paperB}
A.~S.~Petty,
The collision transform,
preprint, 2026.

\bibitem{paper19}
A.~S.~Petty,
The general neutrality theorem,
research note, 2025.

\bibitem{paper20}
A.~S.~Petty,
The double transversality,
research note, 2025.

\bibitem{nfield}
A.~S.~Petty,
\emph{nfield: A structural analysis engine for fractional
field invariants},
software.
\url{https://github.com/alexspetty/nfield}

\end{thebibliography}

\end{document}
